\section{Radar Generation as a Conditional Point Process}
\label{sec:method}

A radar scan contains a variable number of spatial detections, each associated with radar cross section (RCS) and Doppler velocity.
We model this structure directly: a conditional point process describes where returns occur and how many are observed, while spatial mark fields predict their attributes.
Following the point-process perspective of Riedlinger et al.~\cite{riedlinger2026emptyspace}, we derive the learning objective from the observation model.
Our task differs from object detection: we sample radar configurations rather than extract object centers, and use deterministic predictions for the radar marks.

\subsection{Radar Scans as Marked Point Configurations}

Let $C$ denote the camera-derived conditioning representation.
Following RadarGen~\cite{borreda2026radargen}, it comprises aligned BEV appearance, semantic, and radial-velocity maps constructed using estimated depth, semantic segmentation, and optical flow.
The target scan is a finite collection
\begin{equation}
Y=\{(\mathbf{x}_i,\mathbf{m}_i)\}_{i=1}^{N},
\qquad \mathbf{m}_i=(r_i,v_i),
\end{equation}
where $\mathbf{x}_i\in\Omega\subset\mathbb{R}^2$ is a Cartesian BEV position, $r_i$ is RCS in dBsm, and $v_i$ is signed radial velocity in $\mathrm{m}\,\mathrm{s}^{-1}$.
The number of returns $N$ varies between scans and may be zero.
We omit elevation and regard RCS and velocity as marks attached to spatial events.

A neural network maps $C$ to three fields: a log-intensity field and two mark fields.
The intensity defines a distribution over point configurations; the mark fields attach attributes to sampled locations.
Thus, one field prediction supports repeated sampling of distinct radar scans.

\subsection{Conditional Intensity: Count and Location}

We model return locations by an inhomogeneous Poisson point process (PPP) conditioned on $C$.
Its intensity $\lambda_\theta(\mathbf{x}\mid C)$ is nonnegative and measured in expected returns per unit physical area.
For any region $A\subseteq\Omega$,
\begin{equation}
\begin{aligned}
N(A)\mid C &\sim \operatorname{Poisson}(\Lambda_\theta(A\mid C)),\\
\Lambda_\theta(A\mid C)&=\int_A\lambda_\theta(\mathbf{x}\mid C)\,\mathrm{d}\mathbf{x}.
\end{aligned}
\label{eq:regional_count}
\end{equation}
In particular, $\Lambda_\theta=\Lambda_\theta(\Omega\mid C)$ is the expected total number of returns.
Conditional on $N$, locations are independent draws from
\begin{equation}
p_\theta(\mathbf{x}\mid C)
=\frac{\lambda_\theta(\mathbf{x}\mid C)}{\Lambda_\theta}.
\label{eq:normalized_location}
\end{equation}
The intensity therefore encodes two quantities: its integral controls cardinality, while its normalized shape controls the spatial distribution.
A unit-normalized spatial map describes only the latter.

\paragraph{Grid representation.}
We partition $\Omega$ into $K=HW$ equal-area cells.
For each cell $c$, the network predicts $s_c$, with
\begin{equation}
\lambda_{\theta,c}=\frac{\exp(s_c)}{|\Omega|},
\qquad
\mu_c=\lambda_{\theta,c}\frac{|\Omega|}{K}
=\frac{\exp(s_c)}{K}.
\label{eq:cell_mass}
\end{equation}
Here $\mu_c$ is the expected count in the cell and $\Lambda_\theta=\sum_c\mu_c$ for the piecewise-constant intensity.
Equivalently, $s_c$ is a log intensity with respect to spatial measure normalized to unit area.
The factor $1/K$ accounts for cell area; it does not constrain the total predicted count to one.

\subsection{Learning from Observed Returns}

For observed locations $X=\{\mathbf{x}_i\}_{i=1}^N$, the spatial PPP negative log-likelihood, up to terms independent of $\theta$, is
\begin{equation}
\ell_{\mathrm{PPP}}(X,C)
=\Lambda_\theta-\sum_{i=1}^{N}\log\lambda_\theta(\mathbf{x}_i\mid C).
\label{eq:continuous_ppp_nll}
\end{equation}
The integral term penalizes predicted returns throughout the domain, including regions without measurements.
The observation term rewards intensity at measured locations.
For an empty scan, only the integral remains, encouraging low predicted count.

The roles of count and localization become explicit by substituting Eq.~\eqref{eq:normalized_location}:
\begin{equation}
\ell_{\mathrm{PPP}}
=\underbrace{\Lambda_\theta-N\log\Lambda_\theta}_{\text{count}}
-\underbrace{\sum_{i=1}^{N}\log p_\theta(\mathbf{x}_i\mid C)}_{\text{spatial distribution}}.
\label{eq:count_spatial_decomposition}
\end{equation}
The first term is a Poisson count loss up to the constant $\log N!$; the second evaluates the conditional locations.
Both are learned through the same intensity field, without a separate count head.

In the grid implementation, let $c(i)$ be the cell containing return $i$.
We evaluate the equivalent count likelihood
\begin{equation}
\ell_{\mathrm{grid}}
=\sum_c\mu_c-\sum_{i=1}^{N}\log\mu_{c(i)}.
\label{eq:ppp_nll}
\end{equation}
This differs from Eq.~\eqref{eq:continuous_ppp_nll} only by a parameter-independent cell-area term.
Every measured return contributes separately, including returns assigned to the same cell.
We neither binarize cell counts nor interpolate between cells when evaluating the fields.

\paragraph{Mark supervision.}
The network also predicts raw fields $a_c$ and $b_c$.
In the cubic-link configurations considered here, their physical predictions are $\widehat r_c=a_c^3$ and $\widehat v_c=b_c^3$.
These links permit signed, unbounded outputs.
We supervise the predictions directly at measured return locations:
\begin{equation}
\begin{aligned}
\ell_r&=\sum_{i=1}^{N}|\widehat r_{c(i)}-r_i|,\\
\ell_v&=\sum_{i=1}^{N}|\widehat v_{c(i)}-v_i|.
\end{aligned}
\label{eq:mark_losses}
\end{equation}
Errors are computed in physical units, with an ordinary absolute difference for signed velocity.
Multiple returns in one cell retain their individual targets.
Unlike RadarGen's interpolated attribute maps, these losses require no dense mark targets.

L1 supervision admits an interpretation as a fixed-scale Laplace observation model for each mark.
However, our implementation predicts only mark locations: it neither estimates mark scales nor samples residual mark noise.
We therefore use the PPP to model spatial randomness and attach deterministic marks during generation, rather than claim a learned joint density over positions and marks.

\paragraph{Training objective.}
For a single frame, the implemented objective is
\begin{equation}
\mathcal{L}
=\frac{\alpha\ell_{\mathrm{grid}}+\beta\ell_r+\gamma\ell_v}
{\max(N,1)},
\label{eq:full_objective}
\end{equation}
averaged over frames in a minibatch.
We use $\alpha=\beta=\gamma=1$ in the main configuration.
Normalization reduces the contribution of high-count frames relative to unweighted likelihood training.
It yields a count-normalized, likelihood-based objective; it is not the unweighted dataset negative log-likelihood.
For $N=0$, both mark sums vanish and the loss reduces to $\alpha\Lambda_\theta$.

\subsection{Sampling Radar Scans}

A single network forward pass produces all three fields.
We then draw independent cell counts
\begin{equation}
\widetilde n_c\sim\operatorname{Poisson}(\mu_c),
\qquad
\widetilde N=\sum_c\widetilde n_c
\sim\operatorname{Poisson}(\Lambda_\theta).
\end{equation}
For each sampled event in cell $c$, we emit a return at the cell center with marks $(\widehat r_c,\widehat v_c)$.
This is equivalent to sampling the total count first and allocating events to cells with probabilities $\mu_c/\Lambda_\theta$.
Neither an occupancy threshold nor an iterative point-recovery procedure is needed.

Multiple events in a cell are retained and have identical coordinates and marks.
The implemented output is therefore a counting measure that permits repeated returns.
Sampling introduces no subcell jitter or additional mark noise.

\subsection{Field-Prediction Architectures}

The observation model requires three dense output fields and is independent of the network used to predict them.
We investigate two implementations with the same conditioning representation, loss, and sampling procedure.

\paragraph{DiT-based predictor.}
We adapt RadarGen's transformer for direct field prediction.
It processes encoded BEV conditioning at a fixed timestep and uses a trainable decoder to produce the three fields.
The model retains the conditioning autoencoder and a fixed empty-prompt embedding, but replaces iterative denoising with one transformer evaluation followed by PPP sampling.

\paragraph{DINOv3-based predictor.}
A shared pretrained vision transformer processes the three conditioning maps using modality-specific low-rank adapters.
Final-block features are concatenated across modalities and passed to a transformer neck and an attention-based upsampling decoder.
In the main configuration, the backbone, adapters, neck, and decoder are trained jointly.
This implementation requires neither a latent autoencoder nor text or timestep conditioning.

Both receive $512\times512$ conditioning maps.
The DiT model predicts $512\times512$ fields, while the high-resolution DINOv3 model predicts $1024\times1024$ fields over the same physical domain.
Eq.~\eqref{eq:cell_mass} applies at either resolution.

\paragraph{Scope of the observation model.}
The PPP assumes independent counts in disjoint regions conditional on $C$, and its conditional count variance equals its mean.
The network can predict a spatially structured intensity using global scene context, but sampling adds no interactions between returns.
Dependencies caused by shared scatterers, multipath, occlusion, or sensor processing may therefore remain unmodeled.
Grid sampling additionally quantizes locations and assigns the same marks to returns in a cell.
These assumptions provide a tractable statistical model of radar observations rather than a physical model of radar propagation.
